% GENERATED FILE. Edit canonical lessons, then run npm run generate:books.
% canonical-source-hash: fnv1a64:681c6b1665296743
% canonical-lessons: IT-C517-orchestra, IT-C517-sinfonia, IT-C517-spartito, IT-C517-accordare, IT-C517-comporre

\chapter{L'orchestra}
\label{ch:it-l-orchestra}
\hlchaptermodality{517}

\begin{chapteropening}
\textbf{By the end of this chapter:} \emph{I can describe an orchestra and talk about composing music.}

Complete the last lesson of chapter 517: I can describe an orchestra and talk about composing music.
\end{chapteropening}

\section[l'orchestra]{l'orchestra — orchestra}

\label{lesson:IT-C517-orchestra}



\begin{quote}
Before the new one: say the Italian for accordion, then the Italian for to strum.
\end{quote}

\subsection*{You'll want to know: l'orchestra}
\textbf{l'orchestra} — \textquotedblleft{}orchestra\textquotedblright{}.

\begin{quote}
L'orchestra della città suona ogni venerdì.
\end{quote}

\begin{quote}
\emph{The city orchestra plays every Friday.}
\end{quote}

\emph{Il direttore d'orchestra} — the conductor.

\begin{cousinweb}
\emph{orchestra} is a learned borrowing, through Latin, from Greek \emph{orkh\'{\={e}}stra} \textquotedblleft{}the dancing floor of a theatre\textquotedblright{} (EtymDB), from \emph{orkheîsthai} \textquotedblleft{}to dance\textquotedblright{}.
\end{cousinweb}

\begin{grammarlens}[title={What you've built}]
One more word for telling, explaining and describing at B1.
\end{grammarlens}

\subsection*{Your turn}
\begin{itemize}
\raggedright
  \item \emph{Say it:} \emph{l'orchestra}
  \item \emph{Say it:} \emph{l'orchestra}, once more
  \item \emph{Say it:} a sentence of your own with \emph{l'orchestra}
  \item \emph{Recall:} say the Italian for accordion, then the Italian for to strum, then say \emph{l'orchestra} again
\end{itemize}

\begin{tcolorbox}[breakable,colback=teal!4,colframe=teal!35!black,title={Before you move on}]
What does l'orchestra mean? (\textquotedblleft{}Orchestra\textquotedblright{}.) Say it once more.
\end{tcolorbox}
\section[la sinfonia]{la sinfonia — symphony}

\label{lesson:IT-C517-sinfonia}



\begin{quote}
Before the new one: say the Italian for to strum, then the Italian for orchestra.
\end{quote}

\subsection*{You'll want to know: la sinfonia}
\textbf{la sinfonia} — \textquotedblleft{}symphony\textquotedblright{}.

\begin{quote}
Stasera l'orchestra suona una sinfonia famosa.
\end{quote}

\begin{quote}
\emph{Tonight the orchestra is playing a famous symphony.}
\end{quote}

\emph{Sinfonico} — symphonic.

\begin{cousinweb}
\emph{sinfonia} is a learned borrowing from Latin \emph{symphōnia} (EtymDB), from Greek \emph{symphōnía} \textquotedblleft{}sounding together\textquotedblright{}.
\end{cousinweb}

\begin{grammarlens}[title={What you've built}]
One more word for telling, explaining and describing at B1.
\end{grammarlens}

\subsection*{Your turn}
\begin{itemize}
\raggedright
  \item \emph{Say it:} \emph{la sinfonia}
  \item \emph{Say it:} \emph{la sinfonia}, once more
  \item \emph{Say it:} a sentence of your own with \emph{la sinfonia}
  \item \emph{Recall:} say the Italian for to strum, then the Italian for orchestra, then say \emph{la sinfonia} again
\end{itemize}

\begin{tcolorbox}[breakable,colback=teal!4,colframe=teal!35!black,title={Before you move on}]
What does la sinfonia mean? (\textquotedblleft{}Symphony\textquotedblright{}.) Say it once more.
\end{tcolorbox}
\section[lo spartito]{lo spartito — sheet music, score}

\label{lesson:IT-C517-spartito}



\begin{quote}
Before the new one: say the Italian for orchestra, then the Italian for symphony.
\end{quote}

\subsection*{You'll want to know: lo spartito}
\textbf{lo spartito} — \textquotedblleft{}sheet music, score\textquotedblright{}.

\begin{quote}
Ho dimenticato lo spartito a casa.
\end{quote}

\begin{quote}
\emph{I left the sheet music at home.}
\end{quote}

\emph{Leggere lo spartito} — to read music.

\begin{cousinweb}
\emph{spartito} is the past participle of \emph{spartire} \textquotedblleft{}to divide up\textquotedblright{}, built on \emph{partire} \textquotedblleft{}to divide\textquotedblright{}, from Latin \emph{partīrī} (Lewis \& Short): the music is divided into bars.
\end{cousinweb}

\begin{grammarlens}[title={What you've built}]
One more word for telling, explaining and describing at B1.
\end{grammarlens}

\subsection*{Your turn}
\begin{itemize}
\raggedright
  \item \emph{Say it:} \emph{lo spartito}
  \item \emph{Say it:} \emph{lo spartito}, once more
  \item \emph{Say it:} a sentence of your own with \emph{lo spartito}
  \item \emph{Recall:} say the Italian for orchestra, then the Italian for symphony, then say \emph{lo spartito} again
\end{itemize}

\begin{tcolorbox}[breakable,colback=teal!4,colframe=teal!35!black,title={Before you move on}]
What does lo spartito mean? (\textquotedblleft{}Sheet music, score\textquotedblright{}.) Say it once more.
\end{tcolorbox}
\section[accordare]{accordare — to tune (an instrument)}

\label{lesson:IT-C517-accordare}



\begin{quote}
Before the new one: say the Italian for symphony, then the Italian for sheet music.
\end{quote}

\subsection*{You'll want to know: accordare}
\textbf{accordare} — \textquotedblleft{}to tune (an instrument)\textquotedblright{}.

\begin{quote}
Prima del concerto accordo la chitarra.
\end{quote}

\begin{quote}
\emph{Before the concert I tune the guitar.}
\end{quote}

Reflexive \emph{accordarsi} — to agree.

\begin{cousinweb}
\emph{accordare} is inherited from Vulgar Latin \emph{accordāre} (EtymDB), usually linked to Latin \emph{cor, cordis} \textquotedblleft{}heart\textquotedblright{}; the musical sense was helped by \emph{corda} \textquotedblleft{}string\textquotedblright{}.
\end{cousinweb}

\begin{grammarlens}[title={What you've built}]
One more word for telling, explaining and describing at B1.
\end{grammarlens}

\subsection*{Your turn}
\begin{itemize}
\raggedright
  \item \emph{Say it:} \emph{accordare}
  \item \emph{Say it:} \emph{accordare}, once more
  \item \emph{Say it:} a sentence of your own with \emph{accordare}
  \item \emph{Recall:} say the Italian for symphony, then the Italian for sheet music, then say \emph{accordare} again
\end{itemize}

\begin{tcolorbox}[breakable,colback=teal!4,colframe=teal!35!black,title={Before you move on}]
What does accordare mean? (\textquotedblleft{}To tune (an instrument)\textquotedblright{}.) Say it once more.
\end{tcolorbox}
\section[comporre]{comporre — to compose}

\label{lesson:IT-C517-comporre}



\begin{quote}
Before the new one: say the Italian for sheet music, then the Italian for to tune.
\end{quote}

\subsection*{You'll want to know: comporre}
\textbf{comporre} — \textquotedblleft{}to compose\textquotedblright{}.

\begin{quote}
Ha composto una canzone per il compleanno di sua madre.
\end{quote}

\begin{quote}
\emph{He composed a song for his mother's birthday.}
\end{quote}

Like \emph{porre}: \emph{compongo}, \emph{componi}; past participle \emph{composto}. Also to dial: \emph{comporre un numero}.

\begin{cousinweb}
\emph{comporre} is inherited from Latin \emph{compōnere} \textquotedblleft{}to put together, to arrange\textquotedblright{} (EtymDB), from \emph{pōnere} \textquotedblleft{}to put\textquotedblright{}.
\end{cousinweb}

\begin{grammarlens}[title={What you've built}]
One more word for telling, explaining and describing at B1.
\end{grammarlens}

\subsection*{Your turn}
\begin{itemize}
\raggedright
  \item \emph{Say it:} \emph{comporre}
  \item \emph{Say it:} \emph{comporre}, once more
  \item \emph{Say it:} a sentence of your own with \emph{comporre}
  \item \emph{Recall:} say the Italian for sheet music, then the Italian for to tune, then say \emph{comporre} again
\end{itemize}

\begin{tcolorbox}[breakable,colback=teal!4,colframe=teal!35!black,title={Before you move on}]
What does comporre mean? (\textquotedblleft{}To compose\textquotedblright{}.) Say it once more.
\end{tcolorbox}
